\chapter{Limitations and Threats to Validity}\label{ch:limitations}

% --- inlined from paper/sections/06_limitations.tex so the report no longer
% --- depends on the paper's prose, which is cut to the 6-page limit.
\section{Limitations and Threats to Validity}\label{sec:limitations}

We state these before a reader has to raise them.

\textbf{Three datasets, all small and all clinical.} With \NumCells{} cells the
per-cell tests have real power, but generalisation across \emph{datasets} rests on
three points. Nothing here licenses a claim about tabular data in general, and the
difficulty gradient is a design compensation, not a substitute for breadth.

\textbf{Perturbation baselines are off-manifold.} Replacing a feature with the
training median creates rows the data-generating process would not
produce~\cite{hooker2019benchmark}, which can inflate \aopc{}. We report a
marginal-draw variant as a sensitivity check and it did not reverse the ordering, but
both variants share the independence assumption, and a fully on-manifold evaluation
would need a generative model of the data.

\textbf{The plausibility reference sets are not ground truth.} They are compiled from
cited literature and clinical guidelines, frozen before any result existed, and they
remain incomplete, non-causal, and shaped by what their source literature chose to
measure. A low overlap score is as consistent with an incomplete reference set as with
a poor explanation. This is why \RQ{5} is reported as corroborating evidence only.

\textbf{The additivity gate failed for one cell, and we report it rather than widen the
gate.} In PIMA/LightGBM, \AdditivityWorstOutliers{} of 150 explained instances
reconstruct \code{predict\_proba} to worse than the $5\times10^{-3}$ tolerance (worst
\AdditivityWorstMax{}), against a cell median of $1.1\times10^{-9}$, so the
pre-registered rule --- median below tolerance \emph{and} at most $\max(1, 2\%)$ of
instances above it --- is not satisfied and protocol risk R1 fired. The discrepancy is
insensitive to background size (identical at 20 and at all 614 training rows) and we did
not attribute it; the obvious candidate, coincidence with a median-imputed value, was
tested and rejected. It appears to be a library-level edge case at a split boundary in
one fitted model.

We had two pre-registered options: force that cell into log-odds space, which would make
its \shap{} incomparable with \lime{} and so remove it from every cross-explainer
comparison, or widen the gate. We did neither, because widening a gate after seeing the
result it fails is precisely what pre-registration exists to prevent. Instead the cell is
retained, the failure is stated here, and every hypothesis tally is additionally reported
with the cell excluded (\file{results/stats_sensitivity_r1.csv}); no decision in
Table~\ref{tab:hypotheses} changes. A reader who prefers the stricter reading can take
the excluded column.

\textbf{Retraining stability covers three of six models.} Running \RQ{2}c for all six
would have exceeded the compute budget because of \lime{}'s cost, so it is restricted
to the ensembles --- the paper's focus --- and the restriction was registered in
advance rather than discovered afterwards.

\textbf{Timing is single-machine and single-threaded.} The \SpeedFactor{}$\times$ gap
reflects a structural difference in the algorithms and we expect its direction to
hold, but the magnitude is hardware- and implementation-specific.

\textbf{No human evaluation.} Every metric here is a property of the model--explanation
pair. Whether either explanation actually helps a person make a better decision is a
different question, requiring human subjects, and we do not address it.

\textbf{Explanation is not justification.} Following
Rudin~\cite{rudin2019stop}, a faithful post-hoc explanation of a black box is not the
same thing as an interpretable model, and none of these results argue that
post-hoc explanation is sufficient for a high-stakes deployment.

\section{What we would do differently}
Writing this up surfaced choices we would revisit.

\textbf{Widen the dataset axis before widening the model axis.} Six models across three
datasets gave eighteen cells, but the binding constraint on generalisation is the three
datasets, not the six models. Four models across eight datasets would have supported a
stronger claim for the same compute.

\textbf{Choose a plausibility reference differently, or not at all.} \RQ{5} did not
discriminate between explainers and could not have, given that top-5 overlap on an
eight-feature dataset has almost no room to vary. Either use datasets with enough
features for the metric to move, or replace it with a synthetic dataset whose true
feature importances are known by construction --- which would give real ground truth
instead of a literature proxy.

\textbf{Budget for retraining stability from the start.} \RQ{2}c turned out to be the
most interesting stability result and it is the one we had to restrict to three of six
models on cost grounds. Had we known, we would have traded some of \RQ{2}a's re-run
repeats --- which for tree SHAP are provably zero --- for wider retraining coverage.

\textbf{Design the fusion evaluation around complementarity.} With hindsight, testing
fusion between two methods where one is consistently better was testing it under
conditions where it should not help. A more informative design would first identify
cells where the two explainers are close in fidelity but disagree in ranking, and test
fusion there.

\textbf{Keep the amendment log from the first line of code.} We began it only when the
first correction was needed. Three of the amendments in the frozen protocol arose from
tests catching real errors, and recording them from the outset would have made the
sequence clearer.
